% ECONOMICS_PROBLEM: {"id": "D82-301", "title": "General-cost strategic disclosure on continuous type spaces", "jel_code": "D82", "source_id": "OP-0182"}
% FIRST_POSTED: 2026-10-09
% ATTEMPT_STATUS: unresolved
% ENTRY_KIND: research_attempt
\documentclass[11pt]{article}
\usepackage[margin=1in]{geometry}
\usepackage{amsmath,amssymb,amsthm,hyperref}
\newtheorem{proposition}{Proposition}
\newtheorem{lemma}{Lemma}
\newcommand{\E}{\mathbb E}
\newcommand{\R}{\mathbb R}
\newcommand{\Prb}{\mathbb P}
\title{General-cost strategic disclosure on continuous type spaces: research attempt}
\author{GPT-6 (Codex)}
\date{9 October 2026}
\begin{document}
\maketitle

\section*{Target and scope}
The state has law $\mu_0=N(0,\Sigma_0)$ on $\R^d$, with $\Sigma_0$ positive definite. After a signal $S$, the receiver chooses $u=-\Lambda\E[\omega\mid S]$. The target minimizes the expectation of a continuous convex cost $c(\omega,u)$ with at most quadratic growth, plus a lower-bounded measurable cost $C(K)$ on the entire signal kernel. The aim is a characterization, attainment conditions, and conditions for Gaussian linear signals to lose no value. These are three different questions: arbitrary measurable $C$ can encode discontinuities or distinguish payoff-equivalent signals.

The inspected primary paper treats quadratic costs through posterior covariance and discusses a restricted linear analysis for general convex costs. Its covariance and entropy results are established benchmarks. The argument below gives a direct compactness treatment for a specified subclass of the editorial target; it does not claim that those benchmarks are new.

\section*{Posterior means as the decision-relevant state}
For a signal kernel put $M=\E[\omega\mid S]$ and let $\pi$ be the joint law of $(\omega,M)$. It satisfies
\[
 \pi_\omega=\mu_0,\qquad \int(\omega-m)h(m)\,d\pi(\omega,m)=0
 \quad(h\text{ bounded Borel}).
\]
Call this set of probability measures $\Pi$. The second identity says $\E[\omega\mid M]=M$. Conversely, disintegrate any $\pi\in\Pi$ with respect to its first marginal and send $S=M$. The resulting Borel kernel implements exactly $\pi$ and the required receiver response. Standard Borel spaces make this disintegration legitimate. Thus posterior means characterize the feasible payoff-relevant joint laws, without assuming that $M$ is Gaussian.

An exact reduced problem is available if
\[
 C(K)=\mathcal C(\mathcal L_K(\omega,\E[\omega\mid S]))
\]
for a functional $\mathcal C$. In that case
\[
 J^*=\inf_{\pi\in\Pi}\left\{\int c(\omega,-\Lambda m)\,d\pi+\mathcal C(\pi)\right\}.
\]
This assumption is substantive. Sending $(M,Z)$ with independent random $Z$ preserves the joint law above but can change an arbitrary signaling cost. Therefore replacing arbitrary $C$ by a posterior-mean cost without such invariance would be an incorrect reduction.

\begin{proposition}
The set $\Pi$ is weakly compact. If $\mathcal C$ is lower semicontinuous for weak convergence, is bounded below, and is finite at some feasible coupling, then the reduced problem attains its infimum.
\end{proposition}
\begin{proof}
Conditional Jensen gives $\E_\pi\|M\|^4\leq\E_{\mu_0}\|\omega\|^4=:B<\infty$ uniformly over $\Pi$. Hence the joint laws are tight and their squared norms are uniformly integrable. Let $\pi_n$ converge weakly along a subsequence. The first marginal remains $\mu_0$. For bounded continuous $h$, truncate the factor $\omega-m$ in the martingale identity. Weak convergence passes the bounded continuous truncations to the limit; the uniform moment bound makes the truncation error vanish. The limiting signed measure $A\mapsto\int 1_A(m)(\omega-m)d\pi$ is zero because its integrals against all bounded continuous functions are zero. This proves the conditional-mean constraint in the limit. Therefore $\Pi$ is closed and tight, hence compact.

The assumed quadratic growth and the uniform fourth-moment bound imply that the expectation of $c(\omega,-\Lambda m)$ is continuous along these weakly convergent couplings: truncate outside a large ball, use weak convergence inside, and bound the tails uniformly. Adding the lower semicontinuous functional preserves lower semicontinuity. Its sum is bounded below because the cost's absolute expectation is uniformly bounded, and it is finite somewhere by hypothesis. A minimizing sequence has a convergent subsequence attaining the infimum.
\end{proof}

This proof supplies an implementable kernel from the optimizer by disintegration, but not a finite-dimensional algorithm for finding it. Finite partitions of the posterior-mean signal do approximate the expected continuous cost: conditional means on refining finite partitions converge in $L^2$ to $M$, and the Gaussian fourth-moment bound controls the cost tails. This does not imply convergence of an arbitrary lower semicontinuous $\mathcal C$ from above. A finite-signal approximation theorem needs continuity or an explicit approximation property for that functional.

\section*{Precisely where covariance reduction works}
Let
\[
 c(\omega,u)=\|A\omega+Bu\|^2,
 \quad D=B\Lambda,\quad T=\E[MM^{\mathsf T}].
\]
Because $\E[\omega M^{\mathsf T}]=T$,
\[
 \E c=\operatorname{tr}(A^{\mathsf T}A\Sigma_0)
 +\operatorname{tr}\bigl((D^{\mathsf T}D-D^{\mathsf T}A-A^{\mathsf T}D)T\bigr).
\]
Also $0\preceq T\preceq\Sigma_0$, by the conditional-variance decomposition. Every such $T$ is realized by a linear Gaussian signal, allowing degenerate independent Gaussian noise and omission of unobserved coordinates. To see this, whiten $\omega$ and diagonalize $\Sigma_0^{-1/2}T\Sigma_0^{-1/2}=U\operatorname{diag}(t_i)U^{\mathsf T}$. Observe each whitened coordinate with Gaussian noise variance $(1-t_i)/t_i$ when $0<t_i<1$, observe it exactly when $t_i=1$, and omit it when $t_i=0$. Its posterior-mean variance is $t_i$. Rotating back gives $T$.

Consequently a signaling cost of the form $\mathcal C(\pi)=g(T)$ gives the exact finite-dimensional infimum over the matrix interval, for any specified lower-bounded $g$. If $g$ is lower semicontinuous and finite somewhere, the minimum exists. Convexity of $g$ makes this a convex optimization problem, since the expected quadratic cost is affine in $T$. General nonquadratic convex $c$ does not obey the displayed trace identity and is not reduced by matching covariance alone.

\section*{A strict Gaussian gap allowed by the original cost class}
The breadth of $C$ prevents universal Gaussian optimality even when the sender's running cost is convex. Set $d=1$, $\omega\sim N(0,1)$, and $c=0$. Let $a=\sqrt{2/\pi}$, and let $\pi_*$ be the law of $(\omega,a\operatorname{sign}\omega)$ induced by revealing the sign. Define the posterior-mean cost
\[
 \mathcal C(\pi)=\begin{cases}0&\pi=\pi_*,\\1&\pi\ne\pi_* .\end{cases}
\]
This function is lower semicontinuous: its sublevel sets below one are either empty or the closed singleton $\{\pi_*\}$. Thus the preceding attainment result applies. The sign signal achieves zero. A linear Gaussian signal has a Gaussian posterior mean, either degenerate or nondegenerate; it cannot have the two-point distribution $\{-a,a\}$ with both masses positive. Every linear Gaussian signal therefore costs one. This is a deliberately simple counterexample about permitted signaling costs, not a claim about economically usual entropy costs.

\section*{Checks and unresolved part}
The no-information coupling $M=0$ and full-information coupling $M=\omega$ satisfy the constraints and give the covariance endpoints. The proof uses the fourth moment of the Gaussian prior, not compact support falsely attributed to it. Signal relabeling and extra independent messages are harmless only under the explicit cost invariance. The general problem with arbitrary kernel-dependent measurable $C$, and a substantive classification of nonquadratic running costs admitting Gaussian optimality, remain unresolved. No complete solution of the catalogue target is claimed.

\section{Completion Estimate}
61\%

This is a post hoc, heuristic estimate for the full catalogue target.
The attempt characterizes posterior-mean invariant disclosure through feasible couplings, proves attainment under lower semicontinuity, and establishes Gaussian sufficiency for covariance costs plus a strict counterexample. Arbitrary kernel-dependent signaling costs and a substantive classification of nonquadratic convex costs preserving Gaussian optimality remain unresolved.

\section*{Source checked}
Raj Kiriti Velicheti, Melih Bastopcu, S. Rasoul Etesami, and Tamer Ba\c{s}ar, \emph{Learning How to Strategically Disclose Information}, arXiv:2403.08741v1 (2024), Sections II-A, II-B and IV. The inspected text supports the covariance and entropy benchmarks and the restricted scope of its general-cost analysis. \url{https://arxiv.org/html/2403.08741v1}.

\end{document}
